\section{The nonwrapping convolution}
\label{app:nonwrapping}

The main EC construction fixes the coefficient on the oldest of the $m$
feedback symbols to one; the EC literature calls this choice
\emph{wrapping}.  The earlier nonwrapping variant samples that coefficient
with the other feedback coefficients.  Both variants start from an all-zero
feedback state.  We give the nonwrapping enumerator to separate the two
ensembles.

For an output prefix, let state $j\in\{0,\ldots,m\}$ record the number of
trailing zero output bits, capped at $m$.  State $m$ means that all $m$
feedback bits are zero.  From a state $j<m$, the feedback vector is nonzero,
so a uniform binary tap vector makes the next output bit uniform.  From state
$m$, the feedback term is zero and the next output equals the next input.

Let $\mathbf T_m(X,Y)$ be the $(m+1)\times(m+1)$ matrix with nonzero entries
\begin{align}
 (\mathbf T_m)_{j,0}&=\frac{1+X}{2}Y,&
 (\mathbf T_m)_{j,j+1}&=\frac{1+X}{2} &&(0\leq j<m),
 \label{eq:nonwrapping-active}\\
 (\mathbf T_m)_{m,0}&=XY,&
 (\mathbf T_m)_{m,m}&=1.
 \label{eq:nonwrapping-zero}
\end{align}
The exponent of $X$ records input weight and the exponent of $Y$ records
output weight.

\begin{lemma}[Exact nonwrapping enumerator]
\label{lem:nonwrapping-enumerator}
Let $\e_m$ select state $m$.  The expected number of weight-$u$ inputs that
produce a weight-$h$ output under a nonwrapping binary convolution is
\begin{equation}
 [X^uY^h]\,\e_m^{\mathsf T}\mathbf T_m(X,Y)^n\one.
 \label{eq:nonwrapping-enumerator}
\end{equation}
\end{lemma}

\begin{proof}
The two entries from an active feedback state average over the equally likely
zero and one outputs.  The two entries from the zero state copy the input.
Matrix multiplication concatenates the transitions.  The left vector selects
the all-zero initial state, and the right vector sums over the final state.
\end{proof}

For an independent Bernoulli input with nonzero probability $q$, averaging
over the input marker gives a one-variable matrix $\mathbf V_{m,q}(Y)$:
\begin{align}
 (\mathbf V_{m,q})_{j,0}&=Y/2,&
 (\mathbf V_{m,q})_{j,j+1}&=1/2 &&(0\leq j<m),
 \label{eq:nonwrapping-ber-active}\\
 (\mathbf V_{m,q})_{m,0}&=qY,&
 (\mathbf V_{m,q})_{m,m}&=1-q.
 \label{eq:nonwrapping-ber-zero}
\end{align}

\begin{theorem}[Nonwrapping EC first moment]
\label{thm:nonwrapping-first-moment}
Sample $\B\gets\mathcal B^{\mathrm{Ber}}_{k,n,\rho}$ and an independent
nonwrapping convolution $\C$ of memory $m$.  With $a_r$ as in
\eqref{eq:ea-activation},
\begin{equation}
 \Pr_{\B,\C}[\Bad_L(\B\C)]
 \leq
 \sum_{r=1}^{k}\binom kr\sum_{h=0}^{L}[Y^h]\,
 \e_m^{\mathsf T}\mathbf V_{m,a_r}(Y)^n\one.
 \label{eq:nonwrapping-first-moment}
\end{equation}
\end{theorem}

\begin{proof}
A fixed weight-$r$ message has an independent Bernoulli expanded word with
parameter $a_r$.  Equations \eqref{eq:nonwrapping-ber-active} and
\eqref{eq:nonwrapping-ber-zero} give its exact output-weight generating
function.  Summing the low-weight coefficients and then the message shells
gives an expected count of bad messages.  Markov's inequality gives the
claim.
\end{proof}

For any $z\in(0,1]$, the inner coefficient sum in
\eqref{eq:nonwrapping-first-moment} is at most
\begin{equation}
 z^{-L}\e_m^{\mathsf T}\mathbf V_{m,a_r}(z)^n\one.
 \label{eq:nonwrapping-marker}
\end{equation}
This is the same positive-marker relaxation used in the main proof.  The
difference between the wrapping and nonwrapping schemes is therefore confined
to the transfer matrix and its boundary condition.

The finite certificates in this paper target the wrapped construction.  The
matrices above reduce a same-parameter comparison with nonwrapping EC to a
finite numerical calculation, but we make no finite nonwrapping claim.
